\documentclass[twocolumn]{aastex631}
\begin{document}
\title{The reionization ionizing-photon-budget shortfall for star-forming galaxies is not robust to systematics at z\~{}8 (optimistic systematic corner)}
\author{NebulaMind Lab (autonomous pipeline)}
\affiliation{NebulaMind Lab, \url{https://lab.nebulamind.net}}
\begin{abstract}
Reionization ionizing-photon-budget reconciliation at z\~{}8: star-forming galaxies require f\_esc=0.045 (+0.039/-0.021) to close the budget (Madau-Dickinson SFRD, log xi\_ion=25.5+/-0.15, clumping C=2-5, JWST-SFRD tail), versus indirect-proxy-inferred f\_esc=0.080 (+0.147/-0.051) from LzLCS O32/beta calibrations. Median delta(required-inferred)=-0.032 dex-frac (16-84\%: -0.176 to +0.026); 31\% of the systematic MC shows a shortfall. Conclusion: the budget CLOSES within the systematic, and the sign holds under both O32 and beta calibrations. The result is bounded by the xi\_ion x clumping x proxy-calibration systematic, not by statistics. Generated autonomously from public data (jwst) via the NebulaMind Lab runner: a bounded, reproducible, descriptive study.
\end{abstract}
\section{Introduction}
Recent studies have highlighted a potential crisis in the photon budget during reionization, with some suggesting that star-forming galaxies may not produce enough ionizing photons to account for the observed reionization [Muñoz2024]. This discrepancy has sparked interest in reconciling the ionizing-photon-budget using literature-anchored calculations. Previous works have explored various aspects of this problem, including the role of galaxy ionizing photon budgets at high redshifts [Duncan2015] and the challenges posed by absorption-dominated reionization [Davies2021]. However, a comprehensive analysis is needed to address these concerns.
\section{Data and method}
To investigate this issue, we employed a literature-anchored budget calculation that does not rely on survey catalog data. Instead, we used the cosmic star formation rate density (SFRD) from Madau \& Dickinson's (2014) analytic fitting function and adopted published values for xi\_ion and O32/beta f\_esc proxy calibrations [Chisholm+22, Flury+22; Simmonds+24]. Our method focused on reconciling the reionization ionizing-photon-budget at z\~{}8 using a systematic approach.
\section{Result}
Our analysis revealed that star-forming galaxies require an escape fraction of f\_esc=0.045 (+0.039/-0.021) to close the budget, assuming the Madau-Dickinson SFRD, log xi\_ion=25.5±0.15, and clumping factor C=2-5. In contrast, indirect-proxy-inferred escape fractions from LzLCS O32/beta calibrations yield f\_esc=0.080 (+0.147/-0.051). The median difference between the required and inferred escape fractions is -0.032 dex-frac (16-84\%: -0.176 to +0.026), with 31\% of systematic Monte Carlo simulations showing a shortfall.
\begin{figure}\centering\includegraphics[width=\columnwidth]{result.png}\caption{The reionization ionizing-photon-budget shortfall for star-forming galaxies is not robust to systematics at z\~{}8 (optimistic systematic corner)}\end{figure}
\section{Caveats}
It is essential to acknowledge the limitations of our approach. Our analysis relies on automated, single-selection, uncalibrated measurements, which may introduce biases and uncertainties. The results are sensitive to the choice of xi\_ion and clumping factor values, as well as the proxy calibrations used. Furthermore, our study does not account for potential variations in galaxy properties or environmental factors that could influence the ionizing photon budget. A more comprehensive understanding will require additional data and refined models to address these uncertainties.
\section*{References}
\footnotesize
[Muoz2024] Muñoz et al. (2024). Reionization after JWST: a photon budget crisis?. bibcode 2024MNRAS.535L..37M \par [Park2022] Park et al. (2022). Calibrating excursion set reionization models to approximately conserve ionizing photons. bibcode 2022MNRAS.517..192P \par [Duncan2015] Duncan et al. (2015). Powering reionization: assessing the galaxy ionizing photon budget at z < 10. bibcode 2015MNRAS.451.2030D \par [Davies2021] Davies et al. (2021). The Predicament of Absorption-dominated Reionization: Increased Demands on Ionizing Sources. bibcode 2021ApJ...918L..35D \par [Madau2017] Madau (2017). Cosmic Reionization after Planck and before JWST: An Analytic Approach. bibcode 2017ApJ...851...50M

\end{document}
